\documentclass[../main.tex]{subfiles}
\begin{document}
\section{Construction randomness and error semantics}
\label{sec:approx_rel_model}
\begin{definition}[Random approximate relation]
\label{def:approx_rel}
Fix a finite tuple universe $U$ and an exact relation $\Relation{R}\subseteq U$.
A random approximate relation is a random variable with values in
$2^U$, equivalently a family of membership functions
$\ARelation{R}_\theta:U\to\{0,1\}$ indexed by a construction seed $\theta$.
The seed is sampled once; each subsequent probe reads the resulting frozen
representation. For a fixed $t\notin\Relation{R}$ its false-positive rate is
$\Pr_\theta\{\ARelation{R}_\theta(t)=1\}$; for $t\in\Relation{R}$ its
false-negative rate is $\Pr_\theta\{\ARelation{R}_\theta(t)=0\}$.
\end{definition}

Unless stated otherwise, probabilities refer to construction randomness
with the exact relation and query held fixed. The notation
$\ARelation{R}[\fprate][\tprate]$ denotes homogeneous per-tuple rates
$\fprate$ on nonmembers and $\tprate=1-\fnrate$ on members.
These are pointwise construction probabilities, not the probability that
an observed positive is false. Population-average rates require a separate
query distribution; conditional rates on an empty class are undefined.
For a fixed realization, its fraction of errors over a finite candidate set
is a deterministic quantity, varying across constructions.

\subsection{Sources, aliases, and independence}
\label{sec:prob_model}
\begin{assumption}[Independent frozen sources]
\label{asm:source_model}
There is a known source map $\kappa:U\to I$ and fixed exact bits
$x_i\in\{0,1\}$ such that $\Relation{R}(t)=x_{\kappa(t)}$.
The random source bits $X_i$ are mutually independent with
\[
 \Pr(X_i=1)=
 \begin{cases}
   \epsilon_i,&x_i=0,\\
   1-\omega_i,&x_i=1.
 \end{cases}
 \qquad \ARelation{R}_\theta(t)=X_{\kappa(t)}.
\]
Every occurrence of a source reads the same bit. Distinct relations may share
sources, and joint calculations use their combined source map.
\end{assumption}
Aliases must have the same exact truth value for this model. For a symmetric
exact relation, unordered-pair classes satisfy this requirement. Deterministic
tuples, such as forced diagonal members, can be included as degenerate sources.
Independent sources need not imply independent output tuples after a query.

\begin{assumption}[Independent tuple specialization]
\label{asm:ber_trial}
The source map is injective, $\epsilon_i=\fprate$ for exact nonmembers,
and $\omega_i=\fnrate$ for exact members. All distinct tuple outcomes are
mutually independent, including across truth classes.
\end{assumption}
This is a useful idealized channel, not a consequence of having a
Bloom filter~\cite{bf} or a salted hash function. Shared bit arrays, construction
events, and finite hash families can introduce dependence. A concrete
implementation must justify the assumptions it uses; \cref{sec:calibration}
separates an ideal hash model from measurements of an implemented filter.

\paragraph{Shared storage versus shared outcomes.}
The standard perfect-hash-plus-fingerprint construction, also discussed in
the companion filter manuscript~\cite{phf}, illustrates this distinction
\cite{dietzfelbinger2008retrieval}. Place member keys injectively with $g$,
store their $M$-bit fingerprints $V[g(x)]=H(x)$, and accept a key when
$H(x)=V[g(x)]$. Assume $H$ is an ideal random function independent of $g$,
and construction inspects $H$ only on members. Conditional on $g$ and the
stored fingerprints, distinct absent keys have independent uniform
fingerprints. Each match has probability $2^{-M}$, even if the keys share
a slot. Since this probability is constant, independence survives averaging.
Thus shared physical storage need not imply dependent answers; repeated
canonical keys do share an outcome. This idealized realization does not
prove that a particular finite hash implementation meets the same law.

\begin{theorem}[Error counts under independent tuples]
\label{thm:fpbinom}
Under \cref{asm:ber_trial}, for $N$ distinct exact nonmembers the number of
false positives has distribution $\mathrm{Binomial}(N,\fprate)$.
\end{theorem}
\begin{proof}
It is the sum of $N$ independent Bernoulli variables with parameter $\fprate$.
\end{proof}
\begin{theorem}[Realized error fraction]
\label{thm:fpr}
For $N>0$ under the same assumptions, the realized fraction
$\widehat\epsilon_N=\mathrm{FP}_N/N$ has mean $\fprate$ and variance
$\fprate(1-\fprate)/N$. For fixed $0<\fprate<1$, along growing collections
of independent distinct nonmembers,
\[
 \frac{\sqrt N(\widehat\epsilon_N-\fprate)}
 {\sqrt{\fprate(1-\fprate)}}\xrightarrow{d}\mathcal N(0,1).
\]
At $\fprate\in\{0,1\}$ the fraction is deterministic.
\end{theorem}
\begin{proof}
The moment identities follow from the binomial distribution; the standardized
limit is the independent, identically distributed central limit theorem.
\end{proof}
The analogous statement holds for false negatives. With aliases, tuple
counts are weighted sums of source outcomes and need not be binomial.
Repeating one probe $N$ times supplies one source outcome, not $N$ trials.

\subsection{Conditioning and repeated queries}
All later fixed-query guarantees are unconditional over construction, or
conditional on explicitly declared construction information. A query chosen
after observing answers can select a known false positive. Thus a pointwise
construction FPR does not imply the same FPR conditional on a transcript, or
a bound on every query in an adaptive workload; adaptive filters address a
separate problem~\cite{bender2017adaptivity}. For a fixed finite workload, a
union bound on per-query errors remains valid without independence between
queries. Conditioning on a relation-wide property can also destroy the
independent-source law; it must be modeled explicitly, as in conditioning
probabilistic databases~\cite{koch2008conditioning}.

Subset and equality tests remain ordinary Boolean lineages over the joint
source space. For example, two independent approximations agree on one true
tuple with probability $\tprate_1\tprate_2+\fnrate_1\fnrate_2$; the product
over tuples additionally requires tuple independence. The companion set
composition manuscript~\cite{bernoulliComposition} supplies the corresponding
set-level calculations. An approximation compared with itself agrees with
probability one, regardless of its marginal error rates.
\end{document}
